% Figure3.tex — Fig. 3 of the Letter (pgfplots), 510 pt wide, same style as Figs. 1, 2
% (labels \small = 9 pt, tick labels \footnotesize = 8 pt, Computer/Latin Modern).
%   julia --project=.. export_figure3.jl   -> tables/*        lualatex Figure3.tex   -> Figure3.pdf
% Row 1: (a) bifurcation diagram of the m-twisted states with both trajectories. Row 2: (b) enlargement, (c) R and C
% along one cycle of the blue trajectory.
\documentclass[10pt]{standalone}
\usepackage{amsmath,amssymb}
\usepackage{pgfplots}
\pgfplotsset{compat=1.18}
\usetikzlibrary{arrows.meta,plotmarks}
\def\CtMin{1.0593}\def\CtMax{2.0557}

\definecolor{fred}{rgb}{0.64,0.08,0.18}
\definecolor{fblue}{rgb}{0,0.447,0.741}
\definecolor{fbluelight}{rgb}{0.5,0.724,0.871}
\definecolor{fgreen}{rgb}{0.466,0.674,0.188}
\pgfplotsset{
  /pgf/number format/.cd, set thousands separator={}, /pgfplots/.cd,
  every axis/.append style={scaled ticks=false},
  ftfig/.style={scale only axis, font=\footnotesize, label style={font=\small},
    title style={font=\small, at={(0,1)}, anchor=south west, inner sep=0pt, yshift=0.4pt, xshift=-0.5pt},
    major tick length=2pt, tick align=inside, every tick label/.append style={inner sep=1pt},
    xlabel style={inner sep=1pt}, ylabel style={inner sep=1pt}, enlargelimits=false, axis on top, clip mode=individual},
  ftdiag/.style={ftfig, ylabel={$\theta_j'$}, xlabel={$\tau_{ij}/d_{ij}=1/v_j$}},
  mk0/.style={mark=*}, mk1/.style={mark=square*}, mk2/.style={mark=triangle*, mark options={rotate=180}},
  mk3/.style={mark=triangle*}, mk4/.style={mark=diamond*}, mk5/.style={mark=pentagon*},
}
% grey level of branch psi_m (psi_0 black ... psi_5 light grey)
\def\gzero{100}\def\gone{86}\def\gtwo{71}\def\gthree{57}\def\gfour{43}\def\gfive{29}
% one branch: stable solid, unstable dotted, changes of stability as filled markers (#1 = m, #2 = grey level)
\newcommand{\onebranch}[2]{%
  \addplot[black!#2, line width=1.2pt, unbounded coords=jump] table {tables/branch_m#1_stable.dat};
  \addplot[black!#2, line width=1.05pt, densely dotted, unbounded coords=jump] table {tables/branch_m#1_unstable.dat};
  \addplot[only marks, mk#1, mark size=1.7pt, black!#2] table {tables/bif_m#1.dat};}
\newcommand{\allbranches}{\onebranch{5}{\gfive}\onebranch{4}{\gfour}\onebranch{3}{\gthree}\onebranch{2}{\gtwo}\onebranch{1}{\gone}\onebranch{0}{\gzero}}
% red: threshold nu (dotted, from x = #1 to #2), bounds 1/v_max = 0.2 and 1/v_min = 6.8, equilibria of the slow flow
\newcommand{\reds}[2]{%
  \addplot[fred, line width=0.6pt, densely dotted] coordinates {(#1,1) (#2,1)};
  \addplot[fred, line width=0.8pt] coordinates {(0.2,0.85) (0.2,1.15)};
  \addplot[fred, line width=0.8pt] coordinates {(6.8,0.85) (6.8,1.15)};
  \addplot[only marks, mark=o, mark size=1.6pt, fred, line width=0.5pt] table {tables/eq_nu.dat};
  \addplot[only marks, mark=square, mark size=1.5pt, fred, line width=0.5pt] table {tables/eq_bounds.dat};}
% trajectories: blue = switching dynamics (light: node 1, dark: mean), green = frequency selection
\newcommand{\trajs}[1]{% #1 = line width of the light blue band
  \addplot[fbluelight, line width=#1] table {tables/traj_blue_node1.dat};
  \addplot[fblue, line width=0.6pt] table {tables/traj_blue_mean.dat};
  \addplot[fgreen, line width=3.2pt] table {tables/traj_green.dat};}
% red arrow for the slow drift from (#1, #3) to (#2, #3)
\newcommand{\slowarrow}[3]{\draw[fred, line width=0.8pt, -{Stealth[length=4pt]}] (axis cs:#1,#3) -- (axis cs:#2,#3);}
% text label at (#2, #3) with options #1
\newcommand{\lab}[4][]{\node[font=\footnotesize, inner sep=1pt, #1] at (axis cs:#2,#3) {#4};}
\begin{document}
\begin{tikzpicture}
\useasboundingbox (0,0) rectangle (510pt,-232pt);
% ---------------- (a) trajectories first, so that branches and markers stay visible on top of them
\begin{axis}[ftdiag, at={(34pt,-16pt)}, anchor=north west, width=466pt, height=76pt, title={(a)},
    xmin=0, xmax=10, ymin=0.89, ymax=1.1, xtick={0,2,4,6,8,10}, ytick={0.9,1,1.1}, yticklabels={$0.9$,$1.0$,$1.1$}]
  \trajs{3.2pt}\allbranches\reds{0}{10}
  \slowarrow{3.0}{1.6}{1.085}\slowarrow{9.6}{8.2}{1.085}\slowarrow{1.0}{2.4}{0.905}\slowarrow{7.6}{9.0}{0.905}
  \node[fred, anchor=south east, font=\small, inner sep=1pt] at (axis cs:9.95,1.0) {$\nu$};
  \node[fred, anchor=north west, font=\small, inner sep=1pt] at (axis cs:0.25,1.098) {$v_{\max}^{-1}$};
  \node[fred, anchor=north west, font=\small, inner sep=1pt] at (axis cs:6.85,1.098) {$v_{\min}^{-1}$};
  \lab[anchor=south west, black!\gzero]{3.62}{1.05}{$\psi_0$}
  \lab[anchor=south west, black!\gone]{6.22}{0.948}{$\psi_1$}
  \lab[anchor=north, black!\gtwo]{5.3}{0.963}{$\psi_2$}
  \lab[anchor=north, black!\gthree]{1.78}{0.966}{$\psi_3$}
  \lab[anchor=north, black!\gfour]{2.48}{0.953}{$\psi_4$}
  \lab[anchor=north west, black!\gfive]{3.6}{0.909}{$\psi_5$}
  \lab[anchor=west, fblue]{3.47}{1.085}{switching dynamics}
  \lab[anchor=south east, fgreen!80!black]{6.4}{0.891}{frequency selection}
\end{axis}
% ---------------- (b) enlargement, with arrows for the slow drift (red) and the fast jumps (blue)
\begin{axis}[ftdiag, at={(34pt,-134pt)}, anchor=north west, width=210pt, height=76pt, title={(b)},
    xmin=2.92, xmax=3.5, ymin=0.893, ymax=1.08, xtick={3,3.1,3.2,3.3,3.4,3.5}, xticklabels={$3$,$3.1$,$3.2$,$3.3$,$3.4$,$3.5$},
    ytick={0.9,1}, yticklabels={$0.9$,$1.0$}]
  \trajs{1.6pt}\allbranches\reds{2.92}{3.5}
  \addplot[fblue, line width=1pt, -{Stealth[length=4.5pt]}, quiver={u=\thisrow{u}, v=\thisrow{v}}] table[x=x, y=y] {tables/jumps_b.dat};
  \slowarrow{3.47}{3.38}{1.072}\slowarrow{2.95}{3.04}{0.898}
  \node[fred, anchor=north west, font=\small, inner sep=1pt] at (axis cs:2.922,0.999) {$\nu$};
  \lab[anchor=south east, black!\gzero]{3.497}{1.046}{$\psi_0$}
  \lab[anchor=north east, black!\gtwo]{3.497}{1.0285}{$\psi_2$}
  \lab[anchor=south west, black!\gthree]{2.922}{1.0125}{$\psi_3$}
  \lab[anchor=south east, black!\gfour]{3.497}{0.9715}{$\psi_4$}
  \lab[anchor=south east, black!\gfive]{3.497}{0.9135}{$\psi_5$}
\end{axis}
% ---------------- (c) R (black) and C (blue dashed) along one cycle, with the twisted state of each segment
\begin{axis}[ftfig, at={(290pt,-134pt)}, anchor=north west, width=210pt, height=76pt, title={(c)},
    xmin=\CtMin, xmax=\CtMax, ymin=-0.05, ymax=1.08, xtick={1.2,1.4,1.6,1.8,2}, ytick={0,0.5,1},
    xlabel={$t/10^6$}, ylabel={$R$, \textcolor{fblue}{$C$}}]
  \addplot[black, line width=0.9pt] table[x=t, y=R] {tables/panel_c.dat};
  \addplot[fblue, line width=0.9pt, densely dashed] table[x=t, y=C] {tables/panel_c.dat};
  \addplot[only marks, mark=none, nodes near coords={$\psi_{\pgfmathprintnumber[int detect]{\pgfplotspointmeta}}$},
    nodes near coords style={font=\scriptsize, anchor=center, inner sep=0.5pt, fill=white, fill opacity=0.8, text opacity=1},
    point meta=explicit] table[x=t, y=y, meta=m] {tables/labels_c.dat};
  \lab[anchor=south, fblue]{1.225}{0.93}{$C$}
  \lab[anchor=north, black]{1.225}{0.205}{$R$}
\end{axis}
\end{tikzpicture}
\end{document}
